\section{Пространства Соболева. Обобщённые производные}
\label{sec:sobolev-spaces}

\subsection{Обобщённая производная}

Классическая производная требует достаточной гладкости функции. В задачах
математической физики естественно работать с более широким классом функций,
для которых производные существуют только в интегральном, или слабом, смысле.

Пусть $u\in L^1_{\mathrm{loc}}(\Omega)$ и $\alpha=(\alpha_1,\ldots,\alpha_n)$
--- мультииндекс. Функция $v\in L^1_{\mathrm{loc}}(\Omega)$ называется
\emph{обобщённой (слабой) производной} $D^\alpha u$, если
\[
    \int_\Omega u\,D^\alpha\varphi\,dx
    =(-1)^{|\alpha|}\int_\Omega v\,\varphi\,dx
    \qquad
    \forall\varphi\in C_0^\infty(\Omega).
\]
Это определение получается из формулы интегрирования по частям: у пробной
функции $\varphi$ компактный носитель внутри $\Omega$, поэтому граничные члены
исчезают.

Обобщённая производная, если существует, единственна с точностью до множества
меры нуль. Если $u$ имеет обычную классическую производную, то слабая
производная совпадает с ней почти всюду. Важное свойство --- устойчивость при
предельном переходе: если $u_k\to u$ и $D^\alpha u_k\to v$ в подходящих
локальных $L^p$-пространствах, то $v=D^\alpha u$.

\subsection{Определение пространств Соболева}

Для $1\leq p<\infty$ и целого $m\geq0$ пространство Соболева определяется как
\[
    W_p^m(\Omega)=
    \left\{u\in L^p(\Omega):
      D^\alpha u\in L^p(\Omega)
      \ \text{для всех }|\alpha|\leq m\right\}.
\]
Часто используется эквивалентное обозначение $W^{m,p}(\Omega)$. При $m=0$
имеем
\[
    W_p^0(\Omega)=L^p(\Omega).
\]
Для $p=2$ принято обозначение
\[
    H^m(\Omega)=W_2^m(\Omega).
\]

Одна из стандартных норм имеет вид
\[
    \|u\|_{W_p^m(\Omega)}
    =\left(
        \sum_{|\alpha|\leq m}
        \|D^\alpha u\|_{L^p(\Omega)}^p
      \right)^{1/p}.
\]
Пространство $W_p^m(\Omega)$ полно относительно этой нормы, то есть является
банаховым. При $p=2$ пространство $H^m(\Omega)$ является гильбертовым.

\subsection[Пространство Wp0m и плотность гладких функций]{Пространство $W_{p,0}^m$ и плотность гладких функций}

Важное подпространство определяется как замыкание тестовых функций:
\[
    W_{p,0}^m(\Omega)
    =\overline{C_0^\infty(\Omega)}^{\,W_p^m}.
\]
Для $m=1$ оно естественно описывает функции с нулевым граничным значением в
смысле следа.

Гладкие функции позволяют аппроксимировать соболевские функции. Внутри области
это строится с помощью усреднения --- свёртки с гладким ядром. Для
$u\in W_p^m(\Omega)$ и компактного $K\Subset\Omega$ можно построить гладкие
$u_\varepsilon$, для которых
\[
    u_\varepsilon\to u
    \quad\text{в }W_p^m(K).
\]
Чтобы получить глобальное приближение вплоть до границы, используют продолжение
функции за пределы области и разбиение единицы; для этого на границу накладывают
регулярность, например условие Липшица.

\subsection{Неравенство Пуанкаре}

Для ограниченной области $\Omega$ при стандартных условиях регулярности границы
для функций из $W_{p,0}^1(\Omega)$ справедливо \emph{неравенство Пуанкаре}
\[
    \|u\|_{L^p(\Omega)}
    \leq C\|\nabla u\|_{L^p(\Omega)}.
\]
Поэтому на $W_{p,0}^1(\Omega)$ величина $\|\nabla u\|_{L^p}$ задаёт норму,
эквивалентную стандартной норме $W_p^1$. Это утверждение особенно важно в
вариационных постановках краевых задач: оно позволяет получать коэрцитивность
энергетической формы при однородных условиях Дирихле.

\subsection{Теоремы вложения Соболева}

Теоремы вложения показывают, какую дополнительную интегрируемость или гладкость
обеспечивает наличие обобщённых производных. Для ограниченной липшицевой области
$\Omega\subset\mathbb R^n$ типичные результаты для $W_p^1(\Omega)$ имеют вид.

Если $1\leq p<n$, то
\[
    W_p^1(\Omega)\hookrightarrow L^{p^*}(\Omega),
    \qquad
    p^*=\frac{np}{n-p}.
\]
Если $p=n$, то
\[
    W_n^1(\Omega)\hookrightarrow L^q(\Omega)
    \qquad\text{для любого }q<\infty.
\]
Если $p>n$, то элементы $W_p^1(\Omega)$ имеют непрерывных представителей и
вкладываются в пространство Гёльдера:
\[
    W_p^1(\Omega)\hookrightarrow C^{0,\alpha}(\overline\Omega),
    \qquad
    \alpha=1-\frac np.
\]
Смысл этих утверждений: контроль функции и её слабых производных в $L^p$
даёт более сильный контроль самой функции.

Важное усиление --- теорема Реллиха--Кондрашова о компактности вложения. В
частности, при $p<n$ вложение
\[
    W_p^1(\Omega)\hookrightarrow L^q(\Omega)
\]
компактно для $1\leq q<p^*$. Поэтому из ограниченной последовательности в
$W_p^1$ можно выделить подпоследовательность, сильно сходящуюся в $L^q$.

\subsection{След функции на границе}

Для гладкой функции значение $u|_{\partial\Omega}$ определено обычным образом.
Для произвольного элемента $W_p^1(\Omega)$ точечные значения на границе могут
не иметь смысла. Теория следов позволяет корректно продолжить операцию
ограничения на границу.

Для достаточно регулярной, например липшицевой, области существует линейный
непрерывный оператор следа
\[
    \gamma:W_p^1(\Omega)\longrightarrow L^p(\partial\Omega),
\]
совпадающий с обычным ограничением на границу для гладких функций:
\[
    \gamma u=u|_{\partial\Omega},
    \qquad u\in C^1(\overline\Omega).
\]
Имеется оценка
\[
    \|\gamma u\|_{L^p(\partial\Omega)}
    \leq C\|u\|_{W_p^1(\Omega)}.
\]
Именно через оператор следа формулируются граничные условия для слабых решений.

\subsection{Слабая постановка краевой задачи}

Пространства Соболева являются естественной средой для слабых решений
дифференциальных уравнений. Рассмотрим, например, задачу Дирихле для уравнения
Пуассона
\[
    -\Delta u=f\quad\text{в }\Omega,
    \qquad u=0\quad\text{на }\partial\Omega.
\]
Умножим уравнение на тестовую функцию $v$ и проинтегрируем по частям. Получим
\[
    \int_\Omega \nabla u\cdot\nabla v\,dx
    =\int_\Omega f v\,dx.
\]
\emph{Слабая постановка}: найти $u\in H_0^1(\Omega)$ такую, что
\[
    \int_\Omega \nabla u\cdot\nabla v\,dx
    =\int_\Omega f v\,dx
    \qquad\forall v\in H_0^1(\Omega).
\]
Здесь от решения уже не требуется существования классических вторых
производных. Именно такая вариационная форма лежит в основе метода Галёркина и
метода конечных элементов.


\subsection{Примеры слабых производных}

Функция может иметь слабую производную, даже если классическая производная не
существует в отдельных точках. Например,
\[
    u(x)=|x|
\]
принадлежит $W^{1,p}(-1,1)$ для любого конечного $p$, а её слабая производная равна
\[
    u'(x)=\operatorname{sign}x
\]
почти всюду.

В то же время функция Хевисайда имеет распределительную производную, равную дельта-функции,
которая не является функцией из $L^p$; поэтому такая ступенчатая функция не принадлежит
$W^{1,p}$ на интервале, содержащем точку скачка.

\subsection{Банаховость и рефлексивность}

Пространство $W^{m,p}(\Omega)$ с нормой
\[
\|u\|_{W^{m,p}}=
\left(\sum_{|\alpha|\le m}\|D^\alpha u\|_p^p\right)^{1/p}
\]
при $1\le p<\infty$ является банаховым. При $1<p<\infty$ оно рефлексивно.
Для $p=2$ это гильбертово пространство $H^m(\Omega)$ со скалярным произведением,
получаемым суммированием $L^2$-скалярных произведений всех производных до порядка
$m$.

\subsection{Вложения Соболева в типичной форме}

Пусть $\Omega\subset\mathbb R^n$ достаточно регулярна. Для $1\le p<n$ вводят
соболевский показатель
\[
    p^*=\frac{np}{n-p}.
\]
Тогда типичное непрерывное вложение имеет вид
\[
    W^{1,p}(\Omega)\hookrightarrow L^{p^*}(\Omega).
\]
Если $p>n$, функции из $W^{1,p}$ имеют более высокую регулярность: при стандартных
условиях они допускают непрерывных представителей и действует вложение типа Морри
в пространство гёльдеровых функций.

Компактная версия --- теорема Реллиха--Кондрашова: при более слабой целевой норме,
например
\[
    W^{1,p}(\Omega)\Subset L^q(\Omega),\qquad q<p^*,
\]
вложение компактно. Именно компактность используется для извлечения сильно
сходящихся подпоследовательностей.

\subsection{Теорема о следе}

Для липшицевой области существует непрерывный оператор следа, который продолжает
обычное ограничение гладкой функции на границу. В частности, для $H^1(\Omega)$
\[
    \gamma:H^1(\Omega)\to H^{1/2}(\partial\Omega)
\]
непрерывен. Пространство $H_0^1(\Omega)$ можно характеризовать как подпространство
функций $H^1$ с нулевым следом.

\subsection{Лакс--Мильграм для слабой задачи Пуассона}

Для задачи
\[
-\Delta u=f,\qquad u|_{\partial\Omega}=0
\]
слабая постановка:
\[
    \int_\Omega \nabla u\cdot\nabla v\,dx
    =\int_\Omega fv\,dx
    \qquad \forall v\in H_0^1(\Omega).
\]
Билинейная форма
\[
    a(u,v)=\int_\Omega\nabla u\cdot\nabla v\,dx
\]
непрерывна, а по неравенству Пуанкаре коэрцитивна на $H_0^1$. Поэтому по теореме
Лакса--Мильграма при $f$ из подходящего двойственного пространства существует
единственное слабое решение.

\subsection{Что важно сказать на экзамене}

Главное --- дать определение слабой производной через интегрирование по частям
с тестовой функцией и затем определить
\[
    W_p^m(\Omega)=\{u:\ D^\alpha u\in L^p,
    \ |\alpha|\leq m\}.
\]
Нужно отметить полноту $W_p^m$, смысл пространства $W_{p,0}^m$ как замыкания
$C_0^\infty$ и неравенство Пуанкаре, позволяющее контролировать норму функции
через норму её градиента. Затем следует объяснить идею теорем вложения:
производные в $L^p$ дают дополнительную регулярность функции. След позволяет корректно понимать
граничные значения, а слабая постановка переводит краевую задачу в интегральное
тождество в соболевском пространстве.

\newpage
